% Generated from the matching Typst scaffold by
% scripts/migrate_typst_report_to_latex.py.

% ── 5.5  H4: Behavioral Effects of Memory Accumulation ───────────────────
% PURPOSE: Test the analysis-paralysis mechanism — does accumulating memory change HOW
%          the agent acts, and does forgetting fix it? Tie to the GLMM position decline.
% FIGURES/TABLES: behavioral table (tool-calls per policy); GLMM coefficient table.
% CITATIONS: liu2024 (Lost-in-the-Middle); xiong2025 (long-context degradation).
% ─────────────────────────────────────────────────────────────────────────

\section{H4: Behavioral Effects of Memory Accumulation}
\label{sec:h4}

H4 asks a mechanistic question that sits behind the outcome-level results of
Section~\ref{sec:h1}: does carrying more memory change \emph{how} the agent acts, not just
what it costs? The motivating prediction is one of analysis paralysis --- that a fat,
accumulating context makes the agent dither, spending more steps per task, and that
proactive forgetting would relieve this by keeping the working context lean. Because
retrieval is held constant across all six conditions (top-$k=5$; Invariant~\#5), every
policy injects the same \emph{number} of memory items at each step; what differs is only
the quality of the stored pool from which those items are drawn. Any behavioral gap would
therefore trace to accumulation itself, not to a larger injection volume. The result, on
this workload, is that no such gap appears.

\subsection{Step-to-step behavior is flat across policies}

We operationalize behavior as the mean number of tool calls (ReAct steps) the agent issues
per task, averaged over the eight sequences. This is the metric the matrix records
uniformly; we did not collect a separate syntax-error rate, so the behavioral comparison
rests on step frequency alone. Table~\ref{tab:behavioral} reports the per-policy means.

\begin{table}[htbp]
  \centering
  \caption{H4 behavioral metric. Mean tool calls (ReAct steps) per task, averaged over the
  eight sequences, by policy. Values are essentially identical across all six conditions
  (spread $\approx 0.18$ steps), so memory accumulation does not change step-to-step
  behavior. Trustworthy 144-run matrix (\texttt{runs\_144\_seq\_cceb325},
  \texttt{deepseek-v4-flash}).}
  \label{tab:behavioral}
  \begin{tabular}{lr}
    \toprule
    Policy & Mean tool calls / task \\
    \midrule
    cls\_consolidation & 21.11 \\
    full\_memory       & 21.05 \\
    random\_prune      & 21.00 \\
    type\_aware\_decay & 20.98 \\
    no\_memory         & 20.98 \\
    recency\_prune     & 20.93 \\
    \bottomrule
  \end{tabular}
\end{table}

The six values span only $20.93$ to $21.11$ tool calls per task --- a range of about
one-fifth of a single step. \texttt{full\_memory}, the condition that should suffer most
from a bloated context, sits squarely in the middle at $21.05$, and \texttt{no\_memory},
which carries nothing, is indistinguishable from it at $20.98$. The agent expends very
nearly the full $20$-step budget on a typical task regardless of what it remembers; if
anything, the agent is bounded by its step cap rather than by deliberation over a long
context. The predicted analysis-paralysis signature --- more steps under accumulation,
fewer under forgetting --- is simply absent. This is consistent with, and helps explain,
the powered policy-null on resolution in Section~\ref{sec:h1}: if the policies do not act
differently, they have little opportunity to resolve differently.

The mechanism that the long-context literature warns of --- degradation when relevant
content is buried in a swollen window \cite{liu2024,xiong2025} --- is controlled away here
by construction rather than left to chance. Fixing retrieval at top-$k=5$ caps the working
context at a small, constant payload no matter how large the stored pool grows, so
accumulation cannot bloat what the agent actually reads on any step. The flat behavioral
profile is the expected consequence of that design choice: under a bounded injection
budget, the size of the underlying memory store does not reach the agent's reasoning loop.

\subsection{Cost is dominated by the agent loop, not the memory payload}

The same flatness explains why forgetting buys almost nothing on the compute axis
(Section~\ref{sec:h1}). Total token cost per run rises only from $5.74$M for
\texttt{no\_memory} to $6.09$M for \texttt{full\_memory}; the five memory-carrying policies
fall within about $2\%$ of one another ($5.96$--$6.09$M), and the full spread across all
six is roughly $6\%$. With tool-call counts pinned near $21$ for every policy, the
dominant cost term is the ReAct loop itself: each task runs up to twenty steps whose
growing conversation is resent on every step, so the per-task token bill is set by the
trajectory length, not by the comparatively tiny memory payload injected alongside it.
Pruning, accordingly, saves on the footprint axis (the tokens carried per task; see
Section~\ref{sec:h1}) but not on total compute --- a near-null on the cost axis that is the
direct counterpart of the behavioral null reported here.

\subsection{The trajectory effect: performance declines with sequence position}

If memory policy does not move behavior, what does move the outcome within a sequence? We
read this from the task-level GLMM (binomial/logit, crossed random effect on
\texttt{task\_id}; Invariant~\#14), fit over all $4{,}914$ task observations with
\texttt{cls\_consolidation} as the reference policy. Table~\ref{tab:glmm} reports the fixed
effects.

\begin{table}[htbp]
  \centering
  \caption{Task-level GLMM (binomial/logit), $4{,}914$ observations, converged; reference
  policy \texttt{cls\_consolidation}. Coefficients are log-odds with standard errors. All
  five policy contrasts are non-significant ($|\text{coef}| < 2\,\text{SE}$); task
  difficulty and sequence position are the two significant predictors. The fit is
  approximate (\texttt{statsmodels} variational estimation; the canonical R \texttt{lme4}
  fit was not run), so we report signs and the policy-null rather than exact $p$-values; a
  sensitivity check confirms the direction.}
  \label{tab:glmm}
  \begin{tabular}{lrc}
    \toprule
    Fixed effect & Coefficient $\pm$ SE & Significant \\
    \midrule
    full\_memory (vs cls)       & $+0.116 \pm 0.118$ & No \\
    no\_memory (vs cls)         & $-0.133 \pm 0.117$ & No \\
    random\_prune (vs cls)      & $+0.033 \pm 0.117$ & No \\
    recency\_prune (vs cls)     & $-0.092 \pm 0.117$ & No \\
    type\_aware\_decay (vs cls) & $-0.188 \pm 0.117$ & No \\
    \midrule
    task difficulty             & $-2.18 \pm 0.065$  & Yes \\
    sequence position           & $-1.49 \pm 0.086$  & Yes \\
    \bottomrule
  \end{tabular}
\end{table}

The model tells a consistent story with the behavioral table. Every policy coefficient is
smaller than twice its standard error, so the GLMM corroborates the policy-null on outcomes
at the task level: which memory policy is in force does not measurably shift the odds of
resolving a task. The two predictors that do matter are task difficulty
($-2.18 \pm 0.065$) and, of direct relevance to H4, \emph{sequence position}
($-1.49 \pm 0.086$): the agent's effectiveness declines as it moves deeper into a sequence,
and it does so regardless of which policy governs memory. Memory accumulation does not
rescue this decline --- \texttt{full\_memory}'s slightly positive but non-significant
policy coefficient does nothing to offset the strong negative position term.

This is the trajectory counterpart of the flat step-to-step behavior. The within-sequence
degradation is real and substantial in log-odds, but it is a property of running a long
continual-learning trajectory under this agent, not something that any of the six policies
turns on or off. The honest reading of H4 is therefore that its predicted mechanism is
\emph{absent} on this workload: accumulation neither bloats the agent's behavior nor, via
the GLMM, its odds of success, and forgetting has correspondingly little behavioral damage
to repair. The performance decline that does occur is a trajectory-level phenomenon shared
by every condition.
